\section{Gradient Decomposition and Second-Moment Inflation}
\label{app:gradient_proof}

This appendix formalizes how activation bias propagates into gradients and
inflates adaptive optimizer statistics.

\subsection{Notation and setup}

Consider a linear layer with weight matrix $W \in \mathbb{R}^{d \times h}$,
input activation $x_i \in \mathbb{R}^{h}$,
and output gradient $\delta_i \in \mathbb{R}^{d}$.

The per-token gradient is
\[
g_i = \delta_i x_i^\top \in \mathbb{R}^{d \times h}.
\]

Let
\[
G = \frac{1}{N} \sum_{i=1}^N g_i
\]
denote the batch gradient.

Assume activations decompose as
\[
x_i = \mu + s_i,
\qquad
\mathbb{E}[s_i] = 0.
\]

No assumption is made about $\mathbb{E}[\delta_i]$ at this stage.

\subsection{Exact gradient decomposition}

\begin{lemma}[Gradient decomposition identity]
\label{lem:grad_decomp}
The batch gradient satisfies
\[
G
=
\mathbb{E}[\delta]\mu^\top
+
\mathbb{E}[\delta s^\top].
\]
\end{lemma}

\begin{proof}
Substitute $x_i = \mu + s_i$ into
\[
G = \frac{1}{N}\sum_{i=1}^N \delta_i x_i^\top
=
\frac{1}{N}\sum_{i=1}^N \delta_i(\mu^\top + s_i^\top).
\]
Distribute:
\[
G
=
\left(\frac{1}{N}\sum_{i=1}^N \delta_i\right)\mu^\top
+
\frac{1}{N}\sum_{i=1}^N \delta_i s_i^\top.
\]
Taking expectation over batch sampling yields the stated identity.
\end{proof}

\paragraph{Interpretation.}
Even if $\mathbb{E}[\delta]=0$, the second term
$\mathbb{E}[\delta s^\top]$ remains.
However, the first term is a rank-one matrix aligned with $\mu$ and can
produce coherent gradient structure whenever $\mathbb{E}[\delta]\neq 0$.

\subsection{Second-moment structure of gradients}

Adaptive optimizers such as Adam maintain an elementwise second-moment estimate:
\[
v_{ab} = \mathbb{E}[g_{i,ab}^2].
\]

Since
\[
g_{i,ab} = \delta_{i,a} x_{i,b},
\]
we have the exact identity
\[
v_{ab}
=
\mathbb{E}[\delta_{a}^2 x_b^2].
\]

Using $x_b = \mu_b + s_b$, we expand:

\begin{align}
x_b^2
&=
(\mu_b + s_b)^2 \\
&=
\mu_b^2
+
2\mu_b s_b
+
s_b^2.
\end{align}

Therefore,

\begin{equation}
\label{eq:v_expand}
v_{ab}
=
\mu_b^2 \mathbb{E}[\delta_a^2]
+
2\mu_b \mathbb{E}[\delta_a^2 s_b]
+
\mathbb{E}[\delta_a^2 s_b^2].
\end{equation}

\subsection{Bias-induced inflation term}

If $s_b$ is zero-mean and independent of $\delta_a^2$,
then $\mathbb{E}[\delta_a^2 s_b] = 0$, and we obtain:

\begin{equation}
\label{eq:v_clean}
v_{ab}
=
\mu_b^2 \mathbb{E}[\delta_a^2]
+
\mathbb{E}[\delta_a^2 s_b^2].
\end{equation}

The first term is deterministic and shared across tokens.

\begin{theorem}[Bias-driven second-moment inflation]
\label{thm:second_moment_inflation}
Suppose
\[
\mu_b^2 \gg \mathbb{E}[s_b^2]
\]
for a subset of coordinates $b$.
Then for those coordinates,
\[
v_{ab}
\approx
\mu_b^2 \mathbb{E}[\delta_a^2],
\]
and hence the second moment is dominated by the bias term.
\end{theorem}

\begin{proof}
From \eqref{eq:v_clean},
\[
v_{ab}
=
\mu_b^2 \mathbb{E}[\delta_a^2]
+
\mathbb{E}[\delta_a^2 s_b^2].
\]
If $\mu_b^2 \gg \mathbb{E}[s_b^2]$ and $\delta_a^2$ has finite variance,
then
\[
\mu_b^2 \mathbb{E}[\delta_a^2]
\gg
\mathbb{E}[\delta_a^2 s_b^2].
\]
Therefore the first term dominates.
\end{proof}

\paragraph{Consequences.}
Large $\mu_b$ causes certain coordinates to have inflated $v_{ab}$ values.
Since Adam uses effective learning rate proportional to $1/\sqrt{v_{ab}}$,
these coordinates experience suppressed learning rates.

\subsection{Effect of mean removal}

After centering activations,
\[
x'_i = x_i - \mu = s_i.
\]

The new second moment becomes

\[
v'_{ab}
=
\mathbb{E}[\delta_a^2 s_b^2].
\]

\begin{theorem}[Second-moment reduction under mean removal]
\label{thm:v_reduction}
For coordinates where bias dominates,
\[
\frac{v'_{ab}}{v_{ab}}
\approx
\frac{\mathbb{E}[s_b^2]}{\mu_b^2}
\ll 1.
\]
\end{theorem}

\begin{proof}
Divide \eqref{eq:v_clean}:
\[
\frac{v'_{ab}}{v_{ab}}
=
\frac{\mathbb{E}[\delta_a^2 s_b^2]}
{\mu_b^2 \mathbb{E}[\delta_a^2] + \mathbb{E}[\delta_a^2 s_b^2]}.
\]
If $\mu_b^2 \gg \mathbb{E}[s_b^2]$, denominator dominated by first term,
so
\[
\frac{v'_{ab}}{v_{ab}}
\approx
\frac{\mathbb{E}[s_b^2]}
{\mu_b^2}
\ll 1.
\]
\end{proof}

\paragraph{Interpretation.}
Mean removal reduces extreme $v_{ab}$ entries multiplicatively.
This matches empirical reductions of $4\times$ to $50\times$ observed in
Section~2.4B.

\subsection{Summary}

We have shown:

\begin{enumerate}
\item Bias induces a rank-one component in gradients.
\item Bias contributes a deterministic inflation term to second-moment statistics.
\item Mean removal eliminates this term exactly.
\item The reduction factor scales with $\mu_b^2 / \mathbb{E}[s_b^2}$.
\end{enumerate}

This provides a precise mechanism linking activation bias to optimizer
instability.